\apendice{Criterios de convergencia de series}

En este apéndice presentamos los principales criterios de convergencia para series numéricas que son fundamentales en el análisis de funciones de variable compleja.

\subsection{Series de términos positivos}

\begin{teorema}[Criterio de comparación]
    Sean $\sum a_n$ y $\sum b_n$ dos series de términos positivos. Si existe $N \in \mathbb{N}$ tal que $a_n \leq b_n$ para todo $n \geq N$, entonces:
    \begin{enumerate}
        \item Si $\sum b_n$ converge, entonces $\sum a_n$ converge.
        \item Si $\sum a_n$ diverge, entonces $\sum b_n$ diverge.
    \end{enumerate}
\end{teorema}

\begin{teorema}[Criterio de comparación por paso al límite]
    Sean $\sum a_n$ y $\sum b_n$ dos series de términos positivos. Si $\lim_{n \to \infty} \frac{a_n}{b_n} = L$ con $0 < L < \infty$, entonces ambas series tienen el mismo carácter (convergen o divergen simultáneamente).
\end{teorema}

\begin{teorema}[Criterio de la raíz (Cauchy)]
    Sea $\sum a_n$ una serie de términos positivos. Sea $L = \limsup_{n \to \infty} \sqrt[n]{a_n}$. Entonces:
    \begin{enumerate}
        \item Si $L < 1$, la serie converge.
        \item Si $L > 1$, la serie diverge.
        \item Si $L = 1$, el criterio no decide.
    \end{enumerate}
\end{teorema}

\begin{teorema}[Criterio del cociente (D'Alembert)]
    Sea $\sum a_n$ una serie de términos positivos tal que $a_n > 0$ para todo $n$. Si existe $\lim_{n \to \infty} \frac{a_{n+1}}{a_n} = L$, entonces:
    \begin{enumerate}
        \item Si $L < 1$, la serie converge.
        \item Si $L > 1$, la serie diverge.
        \item Si $L = 1$, el criterio no decide.
    \end{enumerate}
\end{teorema}

\begin{teorema}[Criterio de Raabe]
    Sea $\sum a_n$ una serie de términos positivos. Si existe $\lim_{n \to \infty} n\left(1 - \frac{a_{n+1}}{a_n}\right) = L$, entonces:
    \begin{enumerate}
        \item Si $L > 1$, la serie converge.
        \item Si $L < 1$, la serie diverge.
        \item Si $L = 1$, el criterio no decide.
    \end{enumerate}
\end{teorema}

\begin{teorema}[Criterio de condensación de Cauchy]
    Sea $(a_n)$ una sucesión decreciente de términos positivos. Entonces la serie $\sum_{n=1}^\infty a_n$ converge si y solo si la serie $\sum_{k=0}^\infty 2^k a_{2^k}$ converge.
\end{teorema}

\subsection{Series alternadas}

\begin{teorema}[Criterio de Leibniz]
    Sea $(a_n)$ una sucesión de términos positivos decreciente que tiende a cero. Entonces la serie alternada $\sum_{n=1}^\infty (-1)^{n+1} a_n$ converge.
    
    Además, si $S = \sum_{n=1}^\infty (-1)^{n+1} a_n$ y $S_n = \sum_{k=1}^n (-1)^{k+1} a_k$, entonces:
    \[
    |S - S_n| \leq a_{n+1}
    \]
\end{teorema}

\subsection{Convergencia absoluta y condicional}

\begin{definición}[Convergencia absoluta]
    Se dice que la serie $\sum a_n$ converge absolutamente si la serie $\sum |a_n|$ converge.
\end{definición}

\begin{teorema}[Convergencia absoluta implica convergencia]
    Si la serie $\sum a_n$ converge absolutamente, entonces converge.
\end{teorema}

\begin{definición}[Convergencia condicional]
    Se dice que la serie $\sum a_n$ converge condicionalmente si converge pero no converge absolutamente.
\end{definición}

\subsection{Series de potencias}

\begin{definición}[Radio de convergencia]
    Para una serie de potencias $\sum_{n=0}^\infty a_n (z-z_0)^n$, el radio de convergencia $R$ está dado por:
    \[
    \frac{1}{R} = \limsup_{n \to \infty} \sqrt[n]{|a_n|}
    \]
    (Fórmula de Cauchy-Hadamard)
\end{definición}

\begin{teorema}[Fórmula del cociente para el radio de convergencia]
    Si existe $\lim_{n \to \infty} \left|\frac{a_n}{a_{n+1}}\right| = L$, entonces el radio de convergencia es $R = L$.
\end{teorema}

\begin{teorema}[Convergencia de series de potencias]
    Sea $\sum_{n=0}^\infty a_n (z-z_0)^n$ una serie de potencias con radio de convergencia $R > 0$. Entonces:
    \begin{enumerate}
        \item La serie converge absolutamente para $|z-z_0| < R$.
        \item La serie diverge para $|z-z_0| > R$.
        \item Para $|z-z_0| = R$, la convergencia debe estudiarse caso por caso.
    \end{enumerate}
\end{teorema}

\subsection{Criterios específicos para el análisis complejo}

\begin{teorema}[Criterio de Weierstrass]
    Sea $\sum f_n(z)$ una serie de funciones definidas en un conjunto $D \subset \mathbb{C}$. Si existe una serie numérica convergente $\sum M_n$ tal que $|f_n(z)| \leq M_n$ para todo $z \in D$ y todo $n \in \mathbb{N}$, entonces $\sum f_n(z)$ converge uniforme y absolutamente en $D$.
\end{teorema}

\begin{teorema}[Convergencia uniforme de series de potencias]
    Una serie de potencias $\sum_{n=0}^\infty a_n (z-z_0)^n$ con radio de convergencia $R > 0$ converge uniforme y absolutamente en todo disco cerrado $\{z \in \mathbb{C} : |z-z_0| \leq r\}$ con $r < R$.
\end{teorema}

\subsection{Series importantes}

A continuación se presentan algunas series importantes y sus radios de convergencia:

\ejemplo{
\begin{enumerate}
    \item $e^z = \sum_{n=0}^\infty \frac{z^n}{n!}$, $R = \infty$
    
    \item $\sin z = \sum_{n=0}^\infty (-1)^n \frac{z^{2n+1}}{(2n+1)!}$, $R = \infty$
    
    \item $\cos z = \sum_{n=0}^\infty (-1)^n \frac{z^{2n}}{(2n)!}$, $R = \infty$
    
    \item $\frac{1}{1-z} = \sum_{n=0}^\infty z^n$, $R = 1$
    
    \item $\ln(1+z) = \sum_{n=1}^\infty (-1)^{n+1} \frac{z^n}{n}$, $R = 1$
    
    \item $(1+z)^\alpha = \sum_{n=0}^\infty \binom{\alpha}{n} z^n$, $R = 1$ (para $\alpha \notin \mathbb{N}$)
\end{enumerate}
}

\begin{observación}
Los criterios de la raíz y del cociente son especialmente útiles para determinar el radio de convergencia de series de potencias. El criterio de Weierstrass es fundamental para demostrar la convergencia uniforme de series de funciones, lo cual es esencial en el análisis complejo para garantizar propiedades como la continuidad y diferenciabilidad de la función límite.
\end{observación}

\subsection{Formulas importantes}

\begin{teorema}[Formula para series geométrica]
    Sea $|r| < 1$. Entonces:
    \[
    \sum_{n=0}^\infty r^n = \frac{1}{1-r}
    \]
    De forma generica, empezando en el termino $k$:
    \[
    \sum_{n=k}^\infty r^n = \frac{r^k}{1-r}
    \]
\end{teorema}